\documentclass[12pt]{article}
\usepackage{amsmath, amssymb, amsthm, geometry, parskip}
\geometry{a4paper, margin=1in}

\newtheorem{definition}{Definition}
\newtheorem{proposition}{Proposition}
\newtheorem{theorem}{Theorem}
\newtheorem*{remark}{Remark}

\title{On the Geometric Invalidity of Embedding Differencing\\on the Hypersphere for Semantic Change Detection}
\author{}
\date{}

\begin{document}

\maketitle

\section{Setup and Definitions}

Let $f : \mathcal{X} \to \mathbb{R}^d$ be a learned encoder mapping inputs to a $d$-dimensional embedding space. In self-supervised methods such as DINO, the output is $\ell_2$-normalized so that embeddings lie on the unit hypersphere $\mathbb{S}^{d-1} = \{ z \in \mathbb{R}^d : \|z\| = 1 \}$. We consider two observations $x_1, x_2 \in \mathcal{X}$ of the same scene at different times, producing embeddings $z_1 = f(x_1)$ and $z_2 = f(x_2)$.

\begin{definition}[Change Vector]
A \emph{change vector} $\Delta(z_1, z_2)$ is a quantity satisfying:
\begin{enumerate}
    \item[(i)] \textbf{Endpoint-dependence:} $\Delta$ depends only on $z_1$ and $z_2$.
    \item[(ii)] \textbf{Invertibility:} Given $\Delta$ and $z_1$, one can recover $z_2$.
    \item[(iii)] \textbf{Composability:} For any three embeddings $a, b, c$,
    \[
        \Delta(a, b) \oplus \Delta(b, c) = \Delta(a, c)
    \]
    under some group operation $\oplus$.
\end{enumerate}
\end{definition}

\begin{definition}[Semantic Consistency]
A change representation is \emph{semantically consistent} if the same semantic transition (e.g., forest $\to$ building) produces the same change vector regardless of where in the embedding space it occurs. Formally: if $a, a'$ are embeddings of the same semantic class and $b, b'$ are embeddings of another class, then $\Delta(a, b) = \Delta(a', b')$.
\end{definition}

\section{Change Detection in Euclidean Space}

\begin{proposition}
In Euclidean space $\mathbb{R}^d$, the vector difference $\Delta(z_1, z_2) = z_2 - z_1$ satisfies all three properties of Definition~1.
\end{proposition}

\begin{proof}
(i) The difference $z_2 - z_1$ depends only on the endpoints. (ii) Given $\Delta$ and $z_1$, we recover $z_2 = z_1 + \Delta$. (iii) Composability:
\[
    (b - a) + (c - b) = c - a.
\]
The group operation $\oplus$ is vector addition in $(\mathbb{R}^d, +)$.
\end{proof}

\begin{proposition}
In Euclidean space, if cluster centers are fixed points $\{c_1, \dots, c_K\} \subset \mathbb{R}^d$, then change vectors between semantic classes are unique, transitive, and path-independent.
\end{proposition}

\begin{proof}
The change vector from class $i$ to class $j$ is $\Delta_{ij} = c_j - c_i$. This is unique (determined by the two centers), and transitive:
\[
    \Delta_{ij} + \Delta_{jk} = (c_j - c_i) + (c_k - c_j) = c_k - c_i = \Delta_{ik}.
\]
Furthermore, $\Delta_{ij}$ does not depend on any intermediate class or path---it is a function of the endpoints alone.
\end{proof}

\section{Change Detection on the Hypersphere}

On a Riemannian manifold, the natural analog of a difference vector is the logarithmic map. For $z_1, z_2 \in \mathbb{S}^{d-1}$, the logarithmic map at $z_1$ yields the initial velocity of the geodesic from $z_1$ to $z_2$:
\[
    \log_{z_1}(z_2) = \frac{\theta}{\sin\theta}\left(z_2 - \cos\theta \cdot z_1\right),
\]
where $\theta = \arccos\langle z_1, z_2 \rangle$ is the geodesic distance. This vector lives in the tangent space $T_{z_1}\mathbb{S}^{d-1}$, the $(d-1)$-dimensional hyperplane orthogonal to $z_1$.

\begin{proposition}
The logarithmic map on $\mathbb{S}^{d-1}$ does not satisfy the composability property of Definition~1. That is, in general:
\[
    \Gamma_{a \to b}\!\left(\log_a(b)\right) + \log_b(c) \neq \log_a(c),
\]
where $\Gamma_{a \to b}$ denotes parallel transport along the geodesic from $a$ to $b$.
\end{proposition}

\begin{proof}
The failure of composability is a direct consequence of the non-vanishing curvature of $\mathbb{S}^{d-1}$. The unit sphere has constant sectional curvature $\kappa = 1$. By the Ambrose--Singer holonomy theorem, parallel transport around a closed geodesic triangle with vertices $a, b, c$ rotates vectors by an angle equal to the spherical excess:
\[
    \text{Holonomy angle} = A = \alpha + \beta + \gamma - \pi,
\]
where $\alpha, \beta, \gamma$ are the interior angles of the geodesic triangle and $A$ is its area on the unit sphere.

For any non-degenerate triangle, $A > 0$. This means that transporting $\log_a(b)$ to $T_b\mathbb{S}^{d-1}$ and adding $\log_b(c)$ yields a vector that, when transported back to $T_a\mathbb{S}^{d-1}$, differs from $\log_a(c)$ by a rotation of magnitude $A$.

Since DINO embeddings are typically spread across the sphere (large pairwise geodesic distances), the spherical excess is non-negligible, and the composability error is large.
\end{proof}

\section{On Naive Ambient-Space Differencing}

A common heuristic is to compute the ambient Euclidean difference $z_2 - z_1$ directly, ignoring the spherical constraint. We show this is also problematic.

\begin{proposition}
For embeddings on $\mathbb{S}^{d-1}$, the Euclidean difference $z_2 - z_1$ does not preserve semantic consistency (Definition~2). Specifically, $\|z_2 - z_1\|$ depends on the relative orientation of the endpoints, not solely on their semantic content.
\end{proposition}

\begin{proof}
Let $a, a'$ be two embeddings of class A and $b, b'$ be two embeddings of class B, all on $\mathbb{S}^{d-1}$. We have
\[
    \|b - a\|^2 = 2 - 2\langle a, b \rangle, \qquad \|b' - a'\|^2 = 2 - 2\langle a', b' \rangle.
\]
These are equal only if $\langle a, b \rangle = \langle a', b' \rangle$. Since different instances of the same class map to different points on the sphere, the inner product $\langle a, b \rangle$ varies across instances. Moreover, the \emph{direction} of $b - a$ depends on the absolute positions of $a$ and $b$ on the sphere, not on any intrinsic class relationship. The difference vector $b - a$ does not even lie on the sphere---it inhabits a different space than the embeddings themselves.
\end{proof}

\section{On Concatenation-Based Approaches}

An alternative approach avoids differencing entirely, instead concatenating the two embeddings $[z_1; z_2] \in \mathbb{R}^{2d}$ and feeding this to a learned change detection head.

\begin{remark}
Concatenation does not define a change vector in the sense of Definition~1. It sacrifices:
\begin{itemize}
    \item \textbf{Composability:} There is no operation $\oplus$ on concatenated pairs such that $[a; b] \oplus [b; c] = [a; c]$.
    \item \textbf{Generalization:} The learned head must observe every pair type during training; it cannot extrapolate to unseen transition types.
    \item \textbf{Interpretability:} The change ``representation'' is $2d$-dimensional with no geometric meaning.
\end{itemize}
Concatenation is a learned black-box classifier over pairs, not a geometric change representation.
\end{remark}

\section{Resolution: Flat Gaussian Embedding Spaces}

\begin{theorem}
Let $f : \mathcal{X} \to \mathbb{R}^d$ be an encoder whose embedding distribution is regularized toward an isotropic Gaussian (e.g., via variance and decorrelation constraints enforcing $\mathrm{Cov}[f(x)] \approx I$). Then the Euclidean difference $\Delta(z_1, z_2) = z_2 - z_1$ satisfies:
\begin{enumerate}
    \item[(a)] All properties of Definition~1 (endpoint-dependence, invertibility, composability).
    \item[(b)] Semantic consistency (Definition~2), to the extent that the encoder maps same-class inputs to a compact cluster.
    \item[(c)] $\|\Delta\|$ is a calibrated, interpretable change magnitude with consistent scale across the embedding space, since the isotropic Gaussian prior ensures no region of the space is privileged in norm or variance.
\end{enumerate}
\end{theorem}

\begin{proof}
Properties (a) follow from Propositions~1 and~2: Euclidean space is flat (zero curvature), so parallel transport is the identity map, holonomy vanishes, and vector addition is globally defined.

Property (b) follows from the assumption of compact clusters: if $a \approx a' \approx c_A$ and $b \approx b' \approx c_B$ for cluster centers $c_A, c_B$, then
\[
    \Delta(a, b) \approx c_B - c_A \approx \Delta(a', b').
\]

Property (c) follows from the isotropic Gaussian regularization. Since $\mathrm{Cov} \approx I$, all dimensions contribute equally to $\|\Delta\|$, and the expected squared distance between independent draws scales linearly with dimensionality regardless of position in the space. No correction for local curvature, metric distortion, or norm variation is required.
\end{proof}

\section*{Summary}

Embedding differencing for change detection requires a flat geometry. On the hypersphere, change vectors are not composable (Proposition~3), ambient differencing is semantically inconsistent (Proposition~4), and concatenation abandons geometric structure entirely (Section~5). A flat Gaussian embedding space resolves all three issues simultaneously, enabling change detection via simple Euclidean subtraction with calibrated, transitive, path-independent change magnitudes.

\end{document}  